\subsection{Source-instrumented Libgcrypt reproduction of the Curve25519 discrepancy label}
\label{subsec:libgcrypt-cve-2017-0379-source-instrumented}

This experiment is a local, source-instrumented reproduction of the leakage
predicate used in CVE-2017-0379.  It does not target a production key and it
does not claim to be a blind cross-process cache exploit.  Instead, it builds
Libgcrypt~1.7.8, the last 1.7.x release before the CVE-2017-0379 mitigation,
under a private prefix and instruments the Curve25519 Montgomery ladder in
\texttt{mpi/ec.c}.  The instrumentation records the short/long modular-reduction
predicate at the final multiplication of the Montgomery ladder step, i.e., the
line-19 predicate observed by Genkin, Valenta, and Yarom in the historical
Flush+Reload attack.

For a Montgomery curve
\[
  E_A : v^2=u^3+Au^2+u,
\]
the Kummer doubling map is
\[
  D_A(u)=\frac{(u^2-1)^2}{4u(u^2+Au+1)} .
\]
Comparing this with the toric reference law
\[
  T(u)=\frac{2u}{1-u^2}
\]
through the toric subtraction operation gives
\[
  \Delta_A(u)=D_A(u)\ominus T(u)
  =\frac{u^6+5u^4+8Au^3+11u^2-1}
        {2u(u-1)(u+1)(u^2+2Au+3)} .
\]
Thus the discrepancy has poles at \(u=1\) and \(u=-1\).  For Curve25519,
\(A=486662\); the class \(u=1\) is an order-four point on the curve, and
\(u=-1\) is the corresponding order-four class on the quadratic twist.  In
projective Montgomery/Kummer coordinates, \(u=X/Z\), the same exceptional
condition is
\[
  u=\pm1 \quad\Longleftrightarrow\quad X=\pm Z .
\]
The experiment therefore uses the historical attack input \(x=1\) or \(x=-1\)
as a concrete representative of the exceptional divisor of \(\Delta_A\).

The local harness calls Libgcrypt's public low-level elliptic-curve API on a
generated 255-bit laboratory scalar and the chosen order-four Kummer input.  The
instrumented ladder logs one CSV row per scalar-bit iteration.  The important
fields are
\[
  \texttt{short},\quad \texttt{long},\quad
  \texttt{entry\_pole\_plus},\quad \texttt{entry\_pole\_minus},\quad
  \texttt{product\_is\_zero} .
\]
The short/long label is computed exactly as the vulnerable implementation sees
it: the instrumentation replaces the final line-19 call
\[
  \texttt{ec\_mulm(prd->z, prd->z, p1->z, ctx)}
\]
with the equivalent sequence
\[
  \texttt{mpi\_mul(prd->z, prd->z, p1->z)};
  \quad \text{log pre-reduction limb count};
  \quad \texttt{ec\_mod(prd->z, ctx)} .
\]
A reduction is classified as short when the unreduced product has fewer limbs
than the modulus and long otherwise.

Let \(k_s\) denote the scalar bit processed at ladder row \(s\), in the same
most-significant-to-least-significant order used by Libgcrypt's Montgomery
ladder.  After the initial warm-up rows, the expected leakage equation is
\[
  \rho_s = k_s \oplus k_{s-1},
  \qquad
  \rho_s =
  \begin{cases}
    0, & \text{short line-19 reduction},\\
    1, & \text{long line-19 reduction}.
  \end{cases}
\]
The analysis script checks this equation row-by-row and reconstructs the scalar
suffix from the recurrence
\[
  k_s = k_{s-1}\oplus\rho_s,
\]
seeded with the last warm-up bit.  A successful local run produces
\texttt{results/libgcrypt\_1\_7\_8\_analysis.json}; the main success criteria are
\texttt{transition\_rule\_mismatch\_count = 0} and
\texttt{suffix\_recovery\_from\_true\_warmup\_bit.exact = true}.

This experiment supports the following implementation-level interpretation.  The
toric discrepancy pole \(u=\pm1\) is not only an algebraic exceptional set.  In
Libgcrypt~1.7.8 it corresponds to the low-order Kummer classes that force the
Montgomery ladder into states for which the line-19 field product is either
zero/small or full-size.  The variable-size MPI reduction then exposes a binary
short/long arithmetic label, and that label is exactly the scalar-bit transition
label used in the historical Curve25519 side-channel attack.
